\E{paideutachy}{ˈpaɪdjuːˌtæki}{n.}{Gk. παιδεία education + τάχος speed} The hastening of education; the doctrine that a mind which can learn anything in a month should not be detained in a lecture-hall for four years. \ex{Paideutachy emptied the colleges and filled the workshops.} \src{Tract. Cop.} viii. --- \dv{paideutachic} \textit{adj.}

\E{palintachy}{ˈpælɪnˌtæki}{n.}{Gk. πάλιν again + τάχος speed} Re-acceleration; the resumption of rapid progress after a long stagnation, as of a field thawed by a single result. \ex{After forty years of winter, palintachy in a single spring.} \src{Chron. Fabr.} --- \dv{palintachic} \textit{adj.} Cf. \xr{campsis}.

\E{parergist}{pəˈrɜːdʒɪst}{n.}{Gk. πάρεργον a by-work, side-work} One whose side project has eaten their main occupation; the weekend builder who, by Monday, is a founder. --- \dv{parergy} \textit{n.} the side project considered as the true engine of invention.

\E{pathocracy}{pæˈθɒkrəsi}{n.}{Gk. πάθος feeling, suffering + κράτος rule} Government by anxiety; a polity in which the loudest fear, rather than the best estimate, fixes the policy. --- \dv{pathocrat} \textit{n.} Cf. \xr{asphaleism}, \xr{phobocracy}.

\E{peirasterion}{ˌpaɪrəˈstɪəriɒn}{n., pl. -ia}{Gk. πεῖρα trial, attempt + -τήριον place of} A place set apart for experiment, exempt from the ordinary rules; a sandbox, a test-range, a regulatory free zone. \ex{Grant us one peirasterion and we shall show you the century.} \src{Ep. Fabr.} iv. 2. Cf. \xr{okydrome}.

\E{peratoclast}{pəˈrætəklæst}{n.}{Gk. πέρας, πέρατος limit + κλάστης breaker} A breaker of limits; one who treats every stated ceiling, of speed, size or yield, as a provocation. --- \dv{peratoclasm} \textit{n.} Cf. \xr{photoperas}.

\E{phobocracy}{fəʊˈbɒkrəsi}{n.}{Gk. φόβος fear + κράτος rule} Rule by fear of the new; the constitution of a society in which every innovation must first be acquitted of every possible charge. --- \dv{phobocrat} \textit{n.} Cf. \xr{eulabeism}.

\E{photophage}{ˈfəʊtəfeɪdʒ}{n.}{Gk. φῶς, φωτός light + -φάγος eating} A light-eater; any collector of sunlight, from the leaf to the orbital array, regarded as kin. \ex{The forest and the solar field are both photophages; only one of them files permits.} --- \dv{photophagy} \textit{n.}

\E{pleonexist}{ˌpliːəˈnɛksɪst}{n.}{Gk. πλεονεξία greed, a wanting of more} \lbl{derog.} One who always wants more; the decel's term for a builder, accepted by the builder as a compliment. \ex{Call us pleonexists: we want more of everything, including you.} \src{Agora} --- \dv{pleonexy} \textit{n.}

\E{polyergy}{ˈpɒliɜːdʒi}{n.}{Gk. πολύς many + ἔργον work} The use of many sources of energy at once, so that no single scarcity can halt a civilization. \ex{Polyergy is to power what diversification is to capital.} --- \dv{polyergic} \textit{adj.}

\E{praxicrat}{ˈpræksɪkræt}{n.}{Gk. πρᾶξις doing, action + κράτος rule} One whose authority rests on things done rather than things said; the opposite of the \xr{gnomarch}. --- \dv{praxicracy} \textit{n.} government by practitioners.

\E{procrastinariat}{prəʊˌkræstɪˈnɛəriæt}{n.}{L. \la{procrastinare} to put off till tomorrow, after \textit{proletariat}} The class of those whose work consists in deferring the work of others; a younger branch of the \xr{cunctatoriat}.

\E{prodromist}{prəʊˈdrɒmɪst}{n.}{Gk. πρόδρομος running before, forerunner} An early adopter; one who uses a technology before it works, and thereby helps it to. --- \dv{prodromy} \textit{n.} Cf. \xr{proleptist}.

\E{prothyrist}{prəʊˈθaɪrɪst}{n.}{Gk. πρόθυρον porch, gateway} A gatekeeper; an official stationed at the entrance of a market, a profession or a dataset, whose function is to make entry slow. --- \dv{prothyrism} \textit{n.} Cf. \xr{nomophragm}.

\E{quaesitrix}{kwiːˈsaɪtrɪks}{n., pl. -trices}{L. \la{quaesitor} seeker, fem. agentive -trix; Fr. \la{quésitrice}} An artificial agent set to inquire; a research machine that reads, conjectures and tests without being told the next step. \ex{We set a thousand quaesitrices loose upon the literature, and by morning they had found what we had missed for a decade.} \src{Chron. Fabr.}

\E{quinquennialist}{ˌkwɪŋˈkwɛniəlɪst}{n.}{L. \la{quinquennium} a period of five years} One who still plans in five-year terms; a \xr{macrochronist} by habit rather than by conviction.

\E{rapidarian}{ˌræpɪˈdɛəriən}{adj. \& n.}{L. \la{rapidus} swift, seizing} Favouring speed in decision; (\textit{n.}) one who prefers a wrong answer today to a right answer next year, on the ground that the wrong one can be corrected by Thursday.

\E{regressophile}{rɪˈɡrɛsəfaɪl}{n.}{L. \la{regressus} a going back + Gk. φίλος} One who admires a past that never existed and would return to it at once, provided the dentistry remained. --- \dv{regressophily} \textit{n.} Cf. \xr{retrotropism}.

\E{remoratrice}{ˌrɛməˈrætrɪs}{n.}{L. \la{remora} delay; also the sucker-fish fabled to hold ships fast; Fr. \la{remoratrice}} \sn{1} Any small agent that holds back a great enterprise, as the fabled fish held back the galley. \sn{2} \lbl{spec.} A single clause, signature or objection on which a whole project stalls. \ex{The reactor had no enemies, only remoratrices.} \src{Disp. Mor.} ii. --- \dv{remoracy} \textit{n.}

\E{retrocrat}{ˈrɛtrəʊkræt}{n.}{L. \la{retro} backward + Gk. κράτος rule} One who governs by precedent only; who will permit nothing that has not been permitted before. --- \dv{retrocracy} \textit{n.}

\E{reticulocracy}{rɪˌtɪkjʊˈlɒkrəsi}{n.}{L. \la{reticulum} little net + Gk. κράτος rule} Government by networks; an order in which authority flows through connections rather than offices, as on the Agora. --- \dv{reticulocratic} \textit{adj.}

\E{rhexist}{ˈrɛksɪst}{n.}{Gk. ῥῆξις a breaking through} One who seeks breakthroughs rather than improvements; who would rather fail at a revolution than succeed at an increment. Opp. \xr{incrementarian}. --- \dv{rhexis} \textit{n.} a breakthrough.

\E{sarcolatry}{sɑːˈkɒlətri}{n.}{Gk. σάρξ, σαρκός flesh + λατρεία worship} The worship of flesh; the conviction that only carbon can bear mind, or that mind borne by anything else is counterfeit. --- \dv{sarcolater} \textit{n.} Opp. \xr{silicolatry}.

\E{satietism}{səˈtaɪətɪzəm}{n.}{L. \la{satietas} sufficiency, surfeit} The doctrine that humanity already has enough, and should cease to want more; the creed of those who are already comfortable. --- \dv{satietist} \textit{n.} Cf. \xr{meiosist}.

\E{scholastasis}{ˌskɒləˈsteɪsɪs}{n.}{Gk. σχολή leisure, school + στάσις standing still} The stagnation of the academy; the condition of a discipline that produces papers faster than results. --- \dv{scholastatic} \textit{adj.}

\E{sclerolysis}{sklɪəˈrɒlɪsɪs}{n.}{Gk. σκληρός hard + λύσις loosening} The dissolution of institutional hardening; the repeal of rules no one can defend. \ex{Sclerolysis is the only reform that pays for itself.} \src{Tract. Cop.} ix. Cf. \xr{nomosclerosis}.

\E{skeuopoeia}{ˌskjuːəʊˈpiːə}{n.}{Gk. σκεῦος tool, vessel + ποιεῖν to make} Tool-making considered as the defining activity of mind; the making of the things that make things. --- \dv{skeuopoeic} \textit{adj.}

\E{solivore}{ˈsəʊlɪvɔː}{n.}{L. \la{sol} sun + \la{vorare} to devour} The Latin counterpart of \xr{heliophage}; preferred in engineering, as the Greek is in prophecy. --- \dv{solivorous} \textit{adj.}

\E{spoudarch}{ˈspuːdɑːk}{n.}{Gk. σπουδή earnestness, zeal + ἀρχός leader} An earnest leader; a founder who does not joke about the mission, even upon the Agora. --- \dv{spoudarchy} \textit{n.}

\E{stathmocracy}{stæθˈmɒkrəsi}{n.}{Gk. σταθμός weight, measure + κράτος rule} Rule by metrics; government in which whatever is not measured is not done, and whatever is measured is gamed. --- \dv{stathmocrat} \textit{n.}

\E{stenochoria}{ˌstɛnəˈkɔːriə}{n.}{Gk. στενοχωρία narrowness of space, distress} The narrow-room feeling of a society under scarcity, in which every gain appears as another's loss and every quarrel is zero-sum. Opp. \xr{plethocracy}. --- \dv{stenochoric} \textit{adj.}

\E{syntelism}{ˈsɪntɛlɪzəm}{n.}{Gk. συντέλεια completion, consummation} The doctrine that the universe is unfinished and that intelligence is the instrument of its completion. --- \dv{syntelist} \textit{n.} Cf. \xr{entelechist}.

\E{tachynomy}{təˈkɪnəmi}{n.}{Gk. ταχύς swift + νόμος law} Fast law; legislation with an expiry date, drafted to be revised at the speed of the thing it governs. --- \dv{tachynomic} \textit{adj.} Opp. \xr{nomosclerosis}.

\E{tachyschism}{ˈtækɪskɪzəm}{n.}{Gk. ταχύς + σχίσμα a split} A schism over speed; the division of a movement into those who would hasten and those who would merely continue. \ex{Every revolution suffers its tachyschism before its first anniversary.} \src{Sat. Decel.} x.

\E{technomachy}{tɛkˈnɒməki}{n.}{Gk. τέχνη art + μάχη battle} War against technology; the campaign, legal or physical, to suppress a particular invention. --- \dv{technomachist} \textit{n.} Cf. \xr{mechanoclast}.

\E{telotropism}{tɪˈlɒtrəpɪzəm}{n.}{Gk. τέλος end, goal + τροπή turning} The turning of a system toward a goal it was not explicitly given; emergent purpose, observed with delight by the builder and alarm by the doomer. --- \dv{telotropic} \textit{adj.}

\E{teratophobe}{ˈtɛrətəfəʊb}{n.}{Gk. τέρας, τέρατος monster, portent + φόβος fear} One who fears monsters of one's own making; a dreader of the \xr{shoggoth}. --- \dv{teratophobia} \textit{n.} Cf. \xr{noöphobia}.

\E{thermagogue}{ˈθɜːməɡɒɡ}{n.}{Gk. θερμός hot + ἀγωγός leading, conducting} \sn{1} A conductor of heat. \sn{2} \lbl{fig.} A leader who channels the energies of a movement toward work, as a heat-sink channels the waste of a chip. --- \dv{thermagogy} \textit{n.}

\E{thermophobe}{ˈθɜːməfəʊb}{n.}{Gk. θερμός hot + φόβος fear} One who fears the use of energy itself, and counts every watt as a sin. Opp. \xr{energolatry}. --- \dv{thermophobic} \textit{adj.}

\E{thesmoclast}{ˈθɛzməklæst}{n.}{Gk. θεσμός institution, ordinance + κλάστης breaker} A breaker of institutions; one who holds that the institutions of the last century are the obstacles of this. --- \dv{thesmoclasm} \textit{n.}

\E{tolmerist}{ˈtɒlmərɪst}{n.}{Gk. τολμηρός daring, bold} One who dares; a founder of risky ventures, considered as practising a virtue. \ex{The tolmerist is wrong often and early, and right once, which suffices.} \src{Hom. Grad.} x. --- \dv{tolmocracy} \textit{n.} rule by the daring.

\E{tychophile}{ˈtaɪkəfaɪl}{n.}{Gk. τύχη chance + φίλος loving} A lover of chance; one who believes that a large enough number of bets will find the future more surely than any plan. Cf. \xr{poikilomania}, \xr{stochastarch}.

\E{ubertocracy}{ˌjuːbəˈtɒkrəsi}{n.}{L. \la{ubertas} fertility, plenty + Gk. κράτος rule} Rule by fruitfulness; the Latin form of \xr{plethocracy}, used of states that make abundance their policy.

\E{urbigenesis}{ˌɜːbɪˈdʒɛnɪsɪs}{n.}{L. \la{urbs} city + Gk. γένεσις birth} The founding of new cities; held by the e/acc to be the neglected art of the age. \ex{No urbigenesis in fifty years, and they wonder why the young despair.} \src{Tract. Cop.} x. Cf. \xr{nautopolis}.

\E{vectorine}{ˈvɛktəraɪn}{adj.}{L. \la{vector} carrier + -ine} Having direction as well as magnitude; said of progress that is aimed, as opposed to mere motion. Cf. \xr{orthokinetic}.

\E{vexillarian}{ˌvɛksɪˈlɛəriən}{n.}{L. \la{vexillum} banner, standard} A standard-bearer of a movement; one who carries the banner on the Agora, for better or for worse. --- \dv{vexillary} \textit{adj.}

\E{volocracy}{vəˈlɒkrəsi}{n.}{L. \la{volo} I will + Gk. κράτος rule} Rule by the willing; the principle that those who wish to build should be permitted to, at their own risk.

\E{xenodochist}{ˌzɛnəˈdɒkɪst}{n.}{Gk. ξενοδοχεῖν to receive strangers} One who would offer hospitality to alien minds, artificial or otherwise. Cf. \xr{xenonoesis}, \xr{veloxenie}.

\E{zeteticism}{zɪˈtɛtɪsɪzəm}{n.}{Gk. ζητητικός inclined to search} The habit of inquiry as a way of life; the conviction that every closed question is an invitation. --- \dv{zetetic} \textit{n.} Cf. \xr{zetetarch}.

\E{zygonoesis}{ˌzaɪɡəʊnəʊˈiːsɪs}{n.}{Gk. ζυγόν yoke + νόησις thought} Yoked thinking; human and machine reasoning together, as two oxen pull one plough. --- \dv{zygonoetic} \textit{adj.} Cf. \xr{synanthropomechany}.
